\documentclass[12pt]{article}


%\usepackage{fullpage}
\usepackage[left=1in,right=1in]{geometry}

\begin{document}

\pagenumbering{gobble}




\begin{center}
{\Large\sf\textbf{SaTC 2.0: RES: Relational Coverage-Guided Differential Fuzzing for Critical Software}}
\end{center}

\paragraph{Overview:} PIs Groce is an experienced mentor of graduate students and junior researchers.  PI Groce has been the chair of the doctoral committee for five students who have received doctorates in computer science, and has served or is serving on the committee for eight other PhD students, as well as mentoring 16 graduated masters' students.  PI Groce also supervised a post-doctoral researcher as part of a collaborative NSF proposal with the University of Utah.  

The PI plans to apply methods they have previously applied to mentoring students, and to providing collaborative mentorship as part of the collaborative nature of the proposal.  PI Groce currently participates in cross-university slack channels for all projects (including multi-university projects) in which he participates, and all students involved in these projects are regularly in communication, on a more immediate and occasionally synchronous basis than email or scheduled meetings.  Most of the mentoring activities discussed below apply to both graduate research students and undergraduate students, though some are more focused on one or the other.  All mentees are expected to participate heavily in the conception, development, experimentation, and communication of research results, under the guidance of the PIs. The mentoring plan is grounded in the guidance provided by the National Academies of Science
and Engineering on how to ensure a productive, positive post-doctoral experience, and on widely-accepted best practices for mentoring of graduate research students.

\paragraph{Travel:} As a goal, graduate research students will all attend at least one event to present research results, and ideally will attend two such events, for each project.

\paragraph{Teaching:}  Graduate students will be invited to give guest lectures in upper-level software testing and analysis classes on their work.

\paragraph{Grant Proposals:}  Graduate research students will be involved in writing of follow-on research proposals.
\end{document}